\chapter{Posterior Targets and HMC Requirements}
\label{ch:bf-hmc-target-requirements}

HMC samples an unconstrained or transformed parameter vector by integrating
Hamiltonian dynamics for a log target.  In BayesFilter, that log target usually
contains a filtering likelihood.  The target contract must therefore be stated
before discussing sampler tuning.

\begin{definition}[Filtering posterior target]
\label{def:bf-filtering-target}
Let $u \in \R^{d_u}$ denote the coordinates used by the sampler and let
$\theta = T(u)$ be the model parameter transform.  A filtering posterior target
has the form
\[
  \log \pi(u \mid y_{1:T})
  =
  \log p(T(u))
  + \ell(T(u); y_{1:T})
  + \log\left|\det DT(u)\right|,
\]
where $\ell$ is the log likelihood returned by the chosen filtering backend.
\end{definition}

For HMC, the value alone is not enough.  The backend and transform jointly
define the gradient seen by the sampler.  If a custom gradient is used, it must
be the gradient of the same scalar target value, not the gradient of a nearby
surrogate unless the sampler is explicitly corrected for that surrogate.  In
particular, if the filtering backend exposes likelihood derivatives in model
parameter coordinates \(\theta\), the HMC wrapper must still construct the full
sampler-coordinate target in \(u\)-space by adding the prior and transform
Jacobian terms from Definition~\ref{def:bf-filtering-target}.

\section{Finite Failure Policy}
\label{sec:bf-finite-failure-policy}

Invalid parameters are ordinary events during HMC warmup and exploration.  A
BayesFilter target must decide how to handle them.  The preferred policy is to
return a deterministic low log density together with the exact derivative of that
same returned fallback scalar for known invalid regions, while separately
recording diagnostics for the reason.  If the fallback scalar is constant on the
invalid region, the corresponding gradient must be zero.  A finite surrogate
direction that is not the derivative of the returned scalar is not HMC-safe
unless the sampler is explicitly corrected for that surrogate.  A linear-algebra
failure that aborts the process is not an HMC-compatible failure mode.

The policy must be explicit for:
\begin{itemize}
  \item parameter transforms and prior support;
  \item stationarity or determinacy restrictions;
  \item covariance positive-definiteness checks;
  \item failed factorizations or solve residuals;
  \item non-finite likelihood or gradient values.
\end{itemize}

\section{Compilation and Differentiability}
\label{sec:bf-hmc-target-compilation}

When a backend is advertised as JIT or XLA compatible, the claim refers to the
complete value-and-gradient path needed by the sampler.  Partial tracing of an
outer function is not enough if the filter, derivative provider, or failure
branch executes outside the compiled region.

The derivative policy should be one of:
\begin{itemize}
  \item analytic score or Hessian recursion;
  \item custom gradient wrapping a certified analytic derivative;
  \item autodiff through smooth primitive operations;
  \item stopped-gradient approximation used only with an explicit correction or
    diagnostic label;
  \item not HMC-safe.
\end{itemize}
